% RESUMO

{
	\cleardoublepage
	\let\clearpage\relax
	\noindent	SOUSA, M. A. B., \textbf{Implementação de modelo de folga de backlash dinâmico em engrenagens para análise de dinâmica de rotação}. 2026. {\color{red} número\underline{ }de\underline{ }páginas f.} Dissertação de Mestrado, Universidade Federal de Uberlândia, Uberlândia.
	
	\chapter*{\MakeUppercase{Resumo}}
}

Caixas de engrenagens são componentes essenciais em máquinas rotativas industriais, mas fenômenos não lineares como o \textit{backlash} dinâmico e a rigidez de engrenamento variante no tempo (\textit{Time-Varying Mesh Stiffness} - TVMS) introduzem elevada complexidade à sua resposta estrutural global.O objetivo geral deste estudo consistiu em desenvolver um modelo matemático capaz de interpretar de maneira coerente tais efeitos e associá-los à dinâmica de rotação. Para tanto, o modelo não linear de engrenamento de \citeonline{yi2019} foi codificado em Python e integrado à biblioteca \textit{Rotordynamic Open-Source Software} (ROSS), baseada no Método dos Elementos Finitos (MEF). A formulação original de 3 graus de liberdade (\textit{gdl}) por engrenagem foi expandida para 6 \textit{gdl} por nó, viabilizando a inclusão de parâmetros geométricos e forças de engrenagens helicoidais, com base em \citeonline{kubur2004} e \citeonline{mo2025}. Adicionalmente, incorporou-se uma técnica de suavização por tangentes hiperbólicas \cite{walha2006} para lidar descontinuidades numéricas e otimizar a convergência dos integradores temporais. Para a resolução das equações diferenciais de movimento, utilizou-se o método de Newmark combinado com o algoritmo de Newton-Raphson e passo de tempo adaptativo. Um modelo analítico simplificado também foi desenvolvido para determinar as frequências naturais do par engrenado. Os resultados foram validados com a literatura, evidenciando frequências naturais perpendiculares e paralelas à Linha de Ação (\textit{Line of Action} - LOA). Diagramas de bifurcação confirmaram transições claras entre regimes de movimento periódico e caótico em altas velocidades. Por fim, a metodologia integrada foi aplicada a um caso prático real de uma caixa de engrenagens industrial de grande porte operada pela Petrobras (rotação de entrada de 1953 rpm e saída de 12225 rpm) empregando o perfil completo de rigidez de engrenamento calculado pelo ROSS. Para este caso a ferramenta de suavização indicou uma redução de aproximadamente $7.3\%$ no tempo total de processamento computacional sem perda de representatividade física do sistema dinâmico. Conclui-se que a abordagem desenvolvida constitui uma solução robusta e de alta fidelidade para o projeto mecânico e o diagnóstico de falhas em transmissões industriais complexas.

\vspace*{\fill}
\noindent\rule{\textwidth}{1pt}
\textit{Palavras-chave: \textit{Backlash}, Rigidez de Engrenamento, Engrenagem, Análise Dinâmica, ROSS}
